\section*{第 \ref{ch:g10:cells-common-unit} 章——细胞：生命的共同单位}

\begin{solution}{exo:g10:cells-common-unit:1}
每种生物都由一个或多个\omterm{def:g10:cells-common-unit:cell}{细胞}构成；\omterm{def:g10:cells-common-unit:cell}{细胞}是结构和功能的基本单位；
每个\omterm{def:g10:cells-common-unit:cell}{细胞}都来自一个已有的\omterm{def:g10:cells-common-unit:cell}{细胞}，
通过分裂产生。
\end{solution}

\begin{solution}{exo:g10:cells-common-unit:2}
\qty{50}{\micro m}; \qty{2.5}{\micro m}; \qty{0.007}{mm}.
\end{solution}

\begin{solution}{exo:g10:cells-common-unit:3}
$36/600 = \qty{0.06}{mm} = \qty{60}{\micro m}$.
\end{solution}

\begin{solution}{exo:g10:cells-common-unit:4}
仅植物\omterm{def:g10:cells-common-unit:cell}{细胞}具有：\omterm{def:g10:cells-common-unit:organelle}{细胞壁}、\omterm{def:g10:cells-common-unit:organelle}{叶绿体}、大\omterm{def:g10:cells-common-unit:organelle}{液泡}。每个\omterm{def:g10:cells-common-unit:cell}{细胞}都有：
\omterm{def:g10:cells-common-unit:cell}{质膜}（还有\omterm{def:g10:cells-common-unit:cell}{细胞质}和\omterm{def:g10:cells-common-unit:organelle}{核糖体}）。
\end{solution}

\begin{solution}{exo:g10:cells-common-unit:5}
\omterm{def:g10:cells-common-unit:prokaryote}{真核细胞}的 DNA 位于\omterm{def:g10:cells-common-unit:organelle}{细胞核}内，并具有
膜包被的\omterm{def:g10:cells-common-unit:organelle}{细胞器}（如口腔上皮\omterm{def:g10:cells-common-unit:cell}{细胞}、酵母菌）；\omterm{def:g10:cells-common-unit:prokaryote}{原核细胞}
两者都没有，其 DNA 游离于\omterm{def:g10:cells-common-unit:cell}{细胞质}中，大小约为
\qty{1}{\micro m}（如细菌，例如 \emph{E. coli}）。
\end{solution}

\begin{solution}{exo:g10:cells-common-unit:6}
实际长度 $6 \times 5/15 = \qty{2}{\micro m}$。放大倍数：标尺
长 \qty{15}{mm} 对应 \qty{5}{\micro m} $= \qty{0.005}{mm}$，因此
$15/0.005 = 3000$: $\times 3000$.
\end{solution}

\begin{solution}{exo:g10:cells-common-unit:7}
\omterm{def:g10:cells-common-unit:organelle}{核糖体}（\qty{25}{nm}）比光学显微镜的
\omterm{def:g10:cells-common-unit:magnification}{分辨率}（\qty{0.2}{\micro m} $= \qty{200}{nm}$）小十倍：两个相邻的
\omterm{def:g10:cells-common-unit:organelle}{核糖体}，或一个\omterm{def:g10:cells-common-unit:organelle}{核糖体}与其外界，无法区分。
放大倍数再高，也只是把模糊放大。
\end{solution}

\begin{solution}{exo:g10:cells-common-unit:8}
这些颗粒是淀粉，以固体颗粒形式储存在\omterm{def:g10:cells-common-unit:cell}{细胞}内（淀粉
不溶于水）。植物的储备物质储存在\omterm{def:g10:cells-common-unit:cell}{细胞}内部，
而不在\omterm{def:g10:cells-common-unit:cell}{细胞}之间的空隙中。
\end{solution}

\begin{solution}{exo:g10:cells-common-unit:9}
洋葱鳞片生长在地下、黑暗中，并储存养分：
\omterm{def:g10:cells-common-unit:organelle}{叶绿体}在那里毫无用处。苔藓叶生活在光照下，
通过光合作用制造食物，因此其\omterm{def:g10:cells-common-unit:cell}{细胞}中充满\omterm{def:g10:cells-common-unit:organelle}{叶绿体}。
结构相同，配置不同。
\end{solution}

\begin{solution}{exo:g10:cells-common-unit:10}
边长 \qty{1}{\micro m}：表面积 \qty{6}{\micro m^2}，体积
\qty{1}{\micro m^3}，比值为每微米 6。边长 \qty{20}{\micro m}：
表面积 \qty{2400}{\micro m^2}，体积 \qty{8000}{\micro m^3}，比值
0.3。细菌每单位内容物所拥有的膜面积
是大\omterm{def:g10:cells-common-unit:cell}{细胞}的二十倍，且仅靠扩散获取物质；
较大的\omterm{def:g10:cells-common-unit:cell}{细胞}则需要内部运输和\omterm{def:g10:cells-common-unit:organelle}{细胞器}。
\end{solution}

\begin{solution}{exo:g10:cells-common-unit:11}
\qty{2}{h} 共分裂 6 次：$2^6 = 64$ 个\omterm{def:g10:cells-common-unit:cell}{细胞}。\qty{4}{h} 共分裂 12 次：
$2^{12} = 4096$.
\end{solution}

\begin{solution}{exo:g10:cells-common-unit:12}
它具有膜和\omterm{def:g10:cells-common-unit:cell}{细胞质}，并能吸收和排出物质，因此基本符合
\omterm{def:g10:cells-common-unit:cell}{细胞}的定义，但它不能分裂，也不能更新其\omterm{def:g10:chemistry-of-life:families}{蛋白质}：它是一种
为容纳血红蛋白而放弃了部分功能的\omterm{def:g10:cells-common-unit:cell}{细胞}，
由完整的\omterm{def:g10:cells-common-unit:cell}{细胞}产生，并由完整的\omterm{def:g10:cells-common-unit:cell}{细胞}替换。
该定义描述的是一般情形；
红\omterm{def:g10:cells-common-unit:cell}{细胞}则是其中一种特化的末端形式。
\end{solution}

\begin{solution}{exo:g10:cells-common-unit:13}
软木是死亡的组织：胡克看到的是\omterm{def:g10:cells-common-unit:cell}{细胞}的纤维素\omterm{def:g10:cells-common-unit:organelle}{细胞壁}，
其内活的内容物早已消失——只剩空盒子。\omterm{def:g10:cells-common-unit:organelle}{细胞核}只存在于活\omterm{def:g10:cells-common-unit:cell}{细胞}的
\omterm{def:g10:cells-common-unit:cell}{细胞质}中，况且他的显微镜的\omterm{def:g10:cells-common-unit:magnification}{分辨率}，以及缺少染色剂，
也会使其无法被看到。
\end{solution}

\begin{solution}{exo:g10:cells-common-unit:14}
病毒没有\omterm{def:g10:cells-common-unit:cell}{细胞质}，没有自身合成的膜，没有代谢，
也不能分裂，因此不符合\omterm{def:g10:cells-common-unit:cell}{细胞}的定义。依据\omterm{def:g10:cells-common-unit:cell}{细胞}学说，
病毒本身并不是活的；它是一种借用\omterm{def:g10:cells-common-unit:cell}{细胞}的机制来复制自身的颗粒。
是否称其为“活的”取决于定义；
但它不是\omterm{def:g10:cells-common-unit:cell}{细胞}这一点则是确定的。
\end{solution}

\begin{solution}{exo:g10:cells-common-unit:15}
体积 $\pi \times (25\,\unit{\micro m})^2 \times 30000\,\unit{\micro m}
\approx \num{5.9e7}\,\unit{\micro m^3}$，约为立方体
（\qty{8000}{\micro m^3}）体积的 7000 倍。由于纤维呈细长圆柱状，
纤维上任一点与\omterm{def:g10:cells-common-unit:cell}{细胞}膜的距离都不超过 \qty{25}{\micro m}，
因此物质交换仍然迅速；每个\omterm{def:g10:cells-common-unit:organelle}{细胞核}只调控纤维上属于它自己的那一段，
因而没有哪个\omterm{def:g10:cells-common-unit:organelle}{细胞核}需要服务于比普通
\omterm{def:g10:cells-common-unit:cell}{细胞}更大的体积。
\end{solution}

\begin{solution}{pb:g10:cells-common-unit:1}
\textbf{1.} $1.8/9 = \qty{0.2}{mm} = \qty{200}{\micro m}$.

\textbf{2.} $1.8/4 = \qty{0.45}{mm}$；约两个\omterm{def:g10:cells-common-unit:cell}{细胞}可横排其间。

\textbf{3.} $28/0.2 = 140$：图的放大倍数为 $\times 140$。

\textbf{4.} $4/140 \approx \qty{0.029}{mm}$，约为 \qty{30}{\micro m}。

\textbf{5.} \omterm{def:g10:cells-common-unit:organelle}{细胞壁}、\omterm{def:g10:cells-common-unit:cell}{质膜}（紧贴\omterm{def:g10:cells-common-unit:organelle}{细胞壁}）、
\omterm{def:g10:cells-common-unit:cell}{细胞质}、\omterm{def:g10:cells-common-unit:organelle}{细胞核}、\omterm{def:g10:cells-common-unit:organelle}{液泡}。$\times 400$ 下存在但不可见的有：
\omterm{def:g10:cells-common-unit:organelle}{核糖体}（\omterm{def:g10:cells-common-unit:organelle}{线粒体}接近极限）。

\textbf{6.} 它们来自厚度为数层\omterm{def:g10:cells-common-unit:cell}{细胞}、且不断受摩擦的
上皮表面：最外层\omterm{def:g10:cells-common-unit:cell}{细胞}扁平并脱落，
这正是棉签能采集到它们的原因。

\textbf{7.} \qty{0.24}{mm} 略高于肉眼的
\qty{0.1}{mm} 分辨极限，是勉强可见的小点。在 $\times 100$ 下，它占视野的
$0.24/1.8 \approx 13\%$，在 $\times 400$ 下则超过一半；
现有物镜都无法使其恰为视野的四分之一。

\textbf{8.} 细菌。电子显微镜可显示其\omterm{def:g10:cells-common-unit:organelle}{细胞壁}、
膜、\omterm{def:g10:cells-common-unit:organelle}{核糖体}和 DNA 区域，以及没有\omterm{def:g10:cells-common-unit:organelle}{细胞核}。

\textbf{9.} 口腔上皮\omterm{def:g10:cells-common-unit:cell}{细胞}只承担一项工作（保护表面）；
在这一身体中，其他\omterm{def:g10:cells-common-unit:cell}{细胞}为其提供养料、供应氧气并协调其活动；
草履虫则没有其他\omterm{def:g10:cells-common-unit:cell}{细胞}可依靠，
必须作为单个\omterm{def:g10:cells-common-unit:cell}{细胞}完成取食、运动、感受刺激和分裂。

\textbf{10.} 外界现在比\omterm{def:g10:cells-common-unit:cell}{细胞质}更咸，因此水分
穿过\omterm{def:g10:cells-common-unit:cell}{细胞}膜流出\omterm{def:g10:cells-common-unit:cell}{细胞}，\omterm{def:g10:cells-common-unit:cell}{细胞}体积减小并
皱缩。

\textbf{11.} $20^3 = \qty{8000}{\micro m^3} = \num{8e-6}\,\unit{mm^3}$.

\textbf{12.} $6 \times 20^2 = \qty{2400}{\micro m^2}$；比值
$2400/8000 = 0.3$ 每微米。

\textbf{13.} 表面积 \qty{6}{\micro m^2}，体积 \qty{1}{\micro m^3}，
比值 6，较上一题大 20 倍。

\textbf{14.} $\qty{8000}{\micro m^3} = \num{8e-9}\,\unit{cm^3}$，故
约 $\num{8e-9}\,\unit{g} = \qty{8}{ng}$。

\textbf{15.} $\pi \times 3.5^2 \times 2 \approx \qty{77}{\micro m^3}$,
约为立方体体积的百分之一。

\textbf{16.} $\qty{60}{L} = \num{6e16}\,\unit{\micro m^3}$；除以
8000，约 $\num{7.5e12}$ 个\omterm{def:g10:cells-common-unit:cell}{细胞}。

\textbf{17.} 其三分之二，约 $\num{5e12}$。

\textbf{18.} $\num{2.5e13} \times 77 \approx \num{1.9e15}\,\unit{\micro
m^3} \approx \qty{1.9}{L}$。血\omterm{def:g10:cells-common-unit:cell}{细胞}：$0.45 \times 4.5 \approx
\qty{2.0}{L}$——相符。

\textbf{19.} 约 $\num{5e12}$ 个大\omterm{def:g10:cells-common-unit:cell}{细胞}，加上 $\num{2.5e13}$ 个红\omterm{def:g10:cells-common-unit:cell}{细胞}，
总计约 $\num{3e13}$ 个，即数十万亿。
\omterm{def:g10:cells-common-unit:cell}{细胞}大小相差约百倍，且各类\omterm{def:g10:cells-common-unit:cell}{细胞}的比例各不相同，因此任何计数都只是
精度在两三倍以内的估计。

\textbf{20.} 约三十万亿（$\num{3e13}$）个\omterm{def:g10:cells-common-unit:cell}{细胞}，全部
由单个起始\omterm{def:g10:cells-common-unit:cell}{细胞}经连续\omterm{def:g10:cells-common-unit:cell}{细胞}分裂产生，
即受精卵。
\end{solution}
